\section{Geometric fusion and the nine-face relation}
\label{sec:geometric-fusion}

\subsection{The collision space and its actual faces}
\label{gf:sec:gf-space}

The pentagon compares bracketings of four inputs.  For five inputs,
the next relation has nine terms.  Six terms are pentagons, and three
are squares.  These surfaces arise as collision strata of ordered
configurations on the real line.  Their boundary orientations determine
the signs of the relation.

For \(n\geq2\), let
\[
 U_n=\{(x_1,\ldots,x_n):0=x_1<x_2<\cdots<x_n=1\}.
\]
This is the quotient of ordered configurations by translation and
positive dilation.  Its orientation is
\(\omega_n=dx_2\wedge\cdots\wedge dx_{n-1}\).
The space \(U_2\) is a positively oriented point.
For a consecutive block \(S=[a,b]\subseteq[1,n]\), record
\[
 u^S_j=\frac{x_j-x_a}{x_b-x_a},\qquad a\leq j\leq b.
\]
The first and last entries are \(0,1\).
Let \(\overline\Delta_{|S|-2}\) denote the closed simplex of
nondecreasing intermediate entries between these endpoints.
Define \(K_n\) as the closure of
\begin{equation}\label{gf:eq:gf-ratio-closure}
 U_n\longrightarrow
 \prod_{\substack{S\text{ consecutive}\\|S|\geq2}}
        \overline\Delta_{|S|-2},\qquad
 x\longmapsto(u^S)_S.
\end{equation}
The factor \(S=[1,n]\) retains the original normalized positions.
No quotient identifies different orders of the inputs.

When a block collapses, its own normalized coordinates remain in
\eqref{gf:eq:gf-ratio-closure}.  If these coordinates have a further
collision, the ratios for the smaller block retain that configuration.
Thus the boundary records both the collision and its internal shape.
The ratios also recover every ordered three-point ratio, by taking
the block between the outer two indices.  Ratios with another denominator
follow by division and the compact values \(0,\infty\).
Consequently this definition is the ordered one-dimensional
configuration compactification of \cite{GF-Sinha}*{Sections 4.2--4.4}.

A \emph{collision tree} is a rooted planar tree with leaves
\(1,\ldots,n\) in their given order.  Each internal vertex has
at least two children.  An internal edge selects the consecutive
block of leaves below it.  Distinct selected blocks are nested or
disjoint.  The full block corresponds to the root, rather than a
boundary face of the normalized space.
Write \(V(T)\) for the internal vertices and \(E(T)\) for
the internal edges.

\begin{proposition}[Collision strata]\label{gf:prop:gf-strata}
The space \(K_n\) is a compact manifold with corners.
Its face with collision tree \(T\) has interior and closure
\begin{equation}\label{gf:eq:gf-stratum-product}
 F_T^\circ\cong\prod_{v\in V(T)}U_{k_v},\qquad
 F_T\cong\prod_{v\in V(T)}K_{k_v},
\end{equation}
where \(k_v\) is the number of children at \(v\).
The codimension is the number of internal edges.
These faces give a regular cellular decomposition of \(K_n\).
For \(n=5\), its numbers of cells in dimensions \(0,1,2,3\)
are \(14,21,9,1\).
\end{proposition}

\begin{proof}
At the root, group equal limiting positions into maximal consecutive
blocks.  In each nonsingleton block, apply the same rule to its
retained normalized coordinates.  Endpoints \(0,1\) are distinct,
so every vertex has at least two children.  This procedure stops
and gives a unique tree.  The distinct child positions at each
vertex give precisely the factors \(U_{k_v}\).

Conversely, choose these vertex configurations and assign a small
positive scale \(t_e\) to every internal edge.  Insert each child
configuration at the corresponding parent position, multiplied by
its scale.  Iterate from the root and normalize the final endpoints.
As all scales approach zero, every ratio in
\eqref{gf:eq:gf-ratio-closure} has a limit.  For a pair of indices,
their first distinct children at the lowest common ancestor determine
the leading nonzero difference.  Thus the limit recovers the given
vertex configurations and is independent of relative vanishing rates
between disjoint scales.

This also describes the local corner coordinates.  One can invert
the insertion near an open stratum by taking representative child
separations and their ratios to the parent separation.  The separations
at each vertex are nonzero there.  The scales are nonnegative boundary
coordinates.  Allowing further collisions in each vertex configuration
gives exactly the product of closures in
\eqref{gf:eq:gf-stratum-product}.
The manifold and regular-cell assertions agree with
\cite{GF-Sinha}*{Theorems 4.4, 4.15, and 4.19}.

Here is a direct finite description when \(n=5\).
The gaps \(d_j=x_{j+1}-x_j\) form a tetrahedron
\(d_j\geq0\), \(\sum_{j=1}^4d_j=1\).
Replace the vertices for blocks \(1234,2345\) by their inward
directions.  Then replace the proper transforms of the edges for
blocks \(123,234,345\) in the same way.
In local nonnegative coordinates, this replacement writes
\((d_1,\ldots,d_r)=t(v_1,\ldots,v_r)\), with
\(\sum v_j=1\), and retains \(v\) at \(t=0\).
It is the real radial blowup of that corner.
It has the same face decomposition as a sufficiently small truncation.
This follows by replacing the radial coordinate \(t\) with a
strictly increasing coordinate that starts at the truncation radius.
Choose the later truncations inside the earlier coordinate collars.

The resulting facets correspond to the nine proper consecutive
blocks of lengths two, three, and four.  Their intersections
correspond to compatible sets of such blocks.  There are \(21\)
compatible pairs and \(14\) compatible triples.  No four are
compatible.  A pair of incompatible overlapping blocks would require
one leaf to have two different immediate ancestors.
Nested or disjoint blocks give a tree and hence an intersection.
This proves the displayed cell counts and identifies the faces
with those of \eqref{gf:eq:gf-ratio-closure}.
\end{proof}

\subsection{Fusion of chains and the nine-face identity}
\label{gf:sec:gf-chains}

For \(p,q\geq2\), put \(n=p+q-1\), and let \(1\leq i\leq p\).
There is a continuous collision insertion
\begin{equation}\label{gf:eq:gf-gamma}
 \gamma_{p,q,i}:K_p\times K_q\longrightarrow K_n.
\end{equation}
It replaces the \(i\)-th outer point by the infinitesimal inner
configuration.  Its image is the face for
\(S=[i,i+q-1]\).  Formula~\eqref{gf:eq:gf-stratum-product}
makes \(\gamma_{p,q,i}\) a homeomorphism onto that closed face.
At boundary points it retains all internal ratios in both factors.
Thus it specifies a point of the compactification, even though
several ordinary positions coincide.

Write \(C_*(K_n;\mathbb Z)\) for the cellular chain complex of
this face decomposition, with homological grading and differential
\(\partial:C_d\to C_{d-1}\).  Orient each face by the product
of its vertex configurations in root-first, depth-first order.
Cellular cross product and \(\gamma_{p,q,i}\) give the chain map
\begin{equation}\label{gf:eq:gf-chain-map}
 \Gamma_{p,q,i}:C_*(K_p;\mathbb Z)\otimes
 C_*(K_q;\mathbb Z)\longrightarrow C_*(K_n;\mathbb Z),
 \qquad a\otimes b\longmapsto a\circ_i b.
\end{equation}
The tensor differential is
\(\partial(a\otimes b)=\partial a\otimes b+(-1)^{|a|}a\otimes\partial b\).
All chains are finite.  Scalar extension defines the same maps
over any commutative coefficient ring.

\begin{lemma}[Compatibility at actual corners]\label{gf:lem:gf-grafting}
The maps \eqref{gf:eq:gf-chain-map} satisfy nested insertion
associativity.  Interchanging two disjoint inserted chains of
degrees \(r,s\) gives the sign \((-1)^{rs}\).
\end{lemma}

\begin{proof}
At a nested collision, both insertion orders give the same vertex
configurations on the same rooted tree.  Every retained ratio is
determined by these data, so the two continuous maps agree.
For disjoint insertions, exchanging the two factors gives that
same map after the product permutation.  The orientation of their
product changes by \((-1)^{rs}\).

On cellular chains, a product of cells maps homeomorphically to
the corresponding face.  Its oriented boundary is the sum of its
factor boundaries with the tensor signs.  This proves the chain
equation and passes both pointwise insertion identities to chains.
No choice of singular fundamental-chain representatives is needed.
\end{proof}

Let \(m_n\in C_{n-2}(K_n;\mathbb Z)\) be the top cell with
orientation \(\omega_n\).  The following formula fixes every sign:
\begin{theorem}[The oriented collision boundary]\label{gf:thm:gf-boundary}
For \(n\geq3\),
\begin{equation}\label{gf:eq:gf-boundary}
 \partial m_n=
 \sum_{\substack{p,q\geq2\\p+q=n+1}}
 \sum_{i=1}^p
 (-1)^{\,i+(p-i)(q-2)}m_p\circ_i m_q.
\end{equation}
In particular, the first relation after the pentagon is
\begin{align}\label{gf:eq:gf-nine}
 \partial m_5={}&-m_4\circ_1m_2+m_4\circ_2m_2
                 -m_4\circ_3m_2+m_4\circ_4m_2\notag\\
 &-m_3\circ_1m_3-m_3\circ_2m_3-m_3\circ_3m_3\notag\\
 &-m_2\circ_1m_4+m_2\circ_2m_4.
\end{align}
The nine terms have the following collision supports:
\[
\begin{array}{c|c|c}
 S&\text{product of faces}&\text{sign}\\\hline
 12,23,34,45&K_4\times K_2&-,+,-,+\\
 123,234,345&K_3\times K_3&-,-,-\\
 1234,2345&K_2\times K_4&-,+
\end{array}
\]
The boundary of the right side of \eqref{gf:eq:gf-nine} is zero
as an integral cellular one-chain.  Each of the \(21\) edges
occurs twice with opposite signs.
\end{theorem}

\begin{proof}
First suppose \(i<p\).  Use normalized outer coordinates
\(0=y_1<\cdots<y_p=1\) and inner coordinates
\(0=u_1<\cdots<u_q=1\).  A collar of the open collision face is
\[
 x_j=y_j\ (j<i),\qquad
 x_{i+a-1}=y_i+t u_a\ (1\leq a\leq q),\qquad
 x_{j+q-1}=y_j\ (j>i),
\]
where \(t\geq0\) is sufficiently small.
Use domain coordinate order
\((t,y_2,\ldots,y_{p-1},u_2,\ldots,u_{q-1})\).
Subtract the row for the cluster anchor from the other cluster rows.
The pivot for \(t\) is the cluster's last point.
Moving this row first and the inner-coordinate rows last gives
\[
 \omega_n=
 (-1)^{\,i+q-3+(p-i-1)(q-2)}t^{q-2}
 \,dt\wedge\omega_p\wedge\omega_q
 \quad(t>0).
\]
The outward normal at \(t=0\) is \(-\partial_t\).
The additional minus sign gives
\((-1)^{i+(p-i)(q-2)}\).
The vanishing factor \(t^{q-2}\) affects the magnitude but
does not change the orientation for positive \(t\).

For \(i=p\), keep the final endpoint fixed by using
\(x_{p+a-1}=1+t(u_a-1)\).
The first cluster point has derivative \(-1\) with respect to \(t\).
Moving its row past the \(p-2\) outer rows gives determinant
sign \((-1)^{p-1}\).  The outward-normal sign is therefore
\((-1)^p\), which is the same formula at \(i=p\).
Every facet is one of these collisions, so their signed sum proves
\eqref{gf:eq:gf-boundary}.

For \(n=5\), the possibilities are \((p,q)=(4,2),(3,3),(2,4)\).
Substitution gives \eqref{gf:eq:gf-nine} and the table.
At a codimension-two corner, choose its two inward normal coordinates
\(s,t\).  Taking outward boundary twice in the two orders
exchanges \(-\partial_s\) and \(-\partial_t\), giving opposite signs.
Both orders yield the same corner map by Lemma~\ref{gf:lem:gf-grafting}.
These are the only two facets through an open edge of this simple
three-dimensional polytope.  Thus the \(42\) edge incidences cancel
in \(21\) pairs.
\end{proof}

For comparison, the lower formulas in precisely this orientation are
\begin{align*}
 \partial m_3&=-m_2\circ_1m_2+m_2\circ_2m_2,\\
 \partial m_4&=-m_3\circ_1m_2+m_3\circ_2m_2-m_3\circ_3m_2
                  +m_2\circ_1m_3+m_2\circ_2m_3.
\end{align*}
The three squares in \eqref{gf:eq:gf-nine} compare two ternary
homotopies.  They are required by the geometry and cannot be
replaced by additional pentagons.

\begin{example}[One square and its four actual edges]
The face for \(S=234\) has outer points \((1,234,5)\) and
inner points \((2,3,4)\).  It is \(K_3\times K_3\).
Its four edge strata add, respectively, the blocks
\(1234,2345,23,34\).  The first two collide outer points.
The last two collide inner points.  With the product orientation,
\[
 \partial(m_3\circ_2m_3)
   =(\partial m_3)\circ_2m_3-m_3\circ_2(\partial m_3).
\]
The minus sign follows from \(|m_3|=1\).
It is the sign needed for cancellation against the adjacent
pentagonal and square faces in \eqref{gf:eq:gf-nine}.
\end{example}

\subsection{Normal fusion of rational collision coefficients}
\label{gf:sec:gf-rational}

The coefficient map at a collision has a direction different from
\eqref{gf:eq:gf-chain-map}.  Chains enter the collision space through
its face.  Rational functions pull back to its normal neighborhood.
Poles require a formal normal parameter before any leading restriction
can be taken.

Use rational coefficients and unnormalized positions in this section.
Write
\[
 \delta_{ab}=z_b-z_a\quad(a<b),\qquad
 R_n=\mathbb Q[z_1,\ldots,z_n],\qquad
 L_n=R_n[\delta_{ab}^{-1}:a<b].
\]
The positive-distance convention changes the sign of \(z_a-z_b\),
but preserves every pole exponent.  For a tuple
\(N=(N_{ab})_{a<b}\) of nonnegative integers, put
\[
 D_N=\prod_{a<b}\delta_{ab}^{N_{ab}},\qquad
 F_NL_n=D_N^{-1}R_n.
\]
Products send \(F_NL_n\otimes F_ML_n\) into \(F_{N+M}L_n\).

Fix \(S=[i,i+q-1]\) and set \(p=n-q+1\).
Let \(\pi:[1,n]\to[1,p]\) collapse this block to \(i\).
Write \(u_1=0,u_q=1\), and set
\[
 B_S=\mathbb Q[y_1,\ldots,y_p,u_2,\ldots,u_{q-1}]
       [(y_b-y_a)^{-1},(u_d-u_c)^{-1}:a<b,\ c<d].
\]
The substitutions
\begin{equation}\label{gf:eq:gf-one-substitution}
 z_{i+a-1}=y_i+t u_a\quad(1\leq a\leq q),\qquad
 z_j=y_{\pi(j)}\quad(j\notin S)
\end{equation}
give the formal coefficient fusion map
\begin{equation}\label{gf:eq:gf-one-ring}
 \Phi_S:L_n\longrightarrow B_S[[t]][t^{-1}].
\end{equation}
The target consists of formal Laurent series with a finite lower
bound in \(t\).  Its coefficients are rational functions on
the open outer and inner configuration spaces.  It is not the
ring of regular functions on the closed collision face.

\begin{theorem}[Pole order and leading fusion]\label{gf:thm:gf-single-poles}
The map \eqref{gf:eq:gf-one-ring} is an injective algebra map.
For
\(N_S=\sum_{a<b,\ a,b\in S}N_{ab}\), it satisfies
\[
 \Phi_S(F_NL_n)\subseteq t^{-N_S}B_S[[t]].
\]
The leading coefficient map
\begin{equation}\label{gf:eq:gf-leading}
 \sigma_{S,N}:F_NL_n\longrightarrow B_S,
 \qquad f\longmapsto[t^0]\bigl(t^{N_S}\Phi_S(f)\bigr)
\end{equation}
has the explicit value
\begin{equation}\label{gf:eq:gf-leading-formula}
 \sigma_{S,N}(P/D_N)=
 \frac{P(y_{\pi(1)},\ldots,y_{\pi(n)})}
 {\displaystyle
  \prod_{\substack{a<b\\a,b\in S}}(u_{b-i+1}-u_{a-i+1})^{N_{ab}}
  \prod_{\substack{a<b\\\{a,b\}\not\subseteq S}}
             (y_{\pi(b)}-y_{\pi(a)})^{N_{ab}}}.
\end{equation}
The bounds are exact for \(P=1\).
Moreover,
\(\sigma_{S,N+M}(fg)=\sigma_{S,N}(f)\sigma_{S,M}(g)\).
\end{theorem}

\begin{proof}
An internal distance becomes
\(t(u_{b-i+1}-u_{a-i+1})\).
Every other distance becomes \(A+tC\), with
\(A=y_{\pi(b)}-y_{\pi(a)}\) invertible in \(B_S\).
Its negative powers have the unique expansion
\begin{equation}\label{gf:eq:gf-binomial}
 (A+tC)^{-r}
 =A^{-r}\sum_{m\geq0}(-1)^m\binom{r+m-1}{m}(tC/A)^m
 \qquad(r\geq1).
\end{equation}
This proves existence and the stated normal bound.  The polynomial
numerator introduces only nonnegative normal powers.
After multiplication by \(t^{N_S}\), setting \(t=0\)
gives \eqref{gf:eq:gf-leading-formula}.
For \(P=1\), its value is a unit of \(B_S\), so the
normal order cannot improve.  Constant-term evaluation on power
series is multiplicative, which gives the last assertion.

The substitution is birational before expansion:
\(y_i=z_i\), \(t=z_{i+q-1}-z_i\), and
\(u_a=(z_{i+a-1}-z_i)/t\), with the other \(y\)'s unchanged.
It therefore embeds \(L_n\) into the fraction field in \(y,u,t\).
Expansion is injective on the rational functions used here, since
their denominators are powers of \(t\) times formal units.
Thus \(\Phi_S\) is injective.
\end{proof}

The leading maps form a multiplicative family indexed by the chosen
pole bounds.  They are not a single evaluation homomorphism
\(L_n\to B_S\).  For example, an internal inverse maps to
\(t^{-1}(u_b-u_a)^{-1}\), which has no value at \(t=0\).
Choosing a larger bound can also make a previously nonzero leading
coefficient vanish.  The filtration index is part of the map.

\subsection{Nested strata, all pair supports, and the normal domain}
\label{gf:sec:gf-nested}

Fix a collision tree \(T\).  At each vertex \(v\), give its
children normalized coordinates
\(u_{v,1}=0<u_{v,2}<\cdots<u_{v,k_v}=1\).
Let \(t_e\) be the scale of an internal child relative to its parent.
For a leaf \(a\), define its coordinate recursively by
\[
 X_{v,a}=
 \begin{cases}
 u_{v,j},&\text{if the child containing }a\text{ is that leaf},\\
 u_{v,j}+t_eX_{w,a},&\text{if that child is the vertex }w.
 \end{cases}
\]
Set \(z_a=b+\rho X_{\mathrm{root},a}\), where \(b\) is a
translation and \(\rho\) is an outer scale.
These coordinates use the first point of each child as its anchor.
For small positive scales, the points retain their strict order.
The last point is \(b+\rho V(t)\), with \(V(0)=1\).
Thus \(\rho\) differs from the total diameter by a positive unit.

For clarity, put
\[
 Q_T=\mathbb Q[b,\rho,(u_{v,j})_{1<j<k_v}]
 \bigl[(u_{v,d}-u_{v,c})^{-1}:v,\ c<d\bigr],
\]
\begin{equation}\label{gf:eq:gf-domain}
 \mathscr L_T=
 Q_T[[t_e:e\in E(T)]][\rho^{-1},t_e^{-1}:e\in E(T)].
\end{equation}
There is one common finite lower bound for each normal exponent.
Formal coefficients may have different orders of tangential poles.
For a corolla, the product of normal variables is empty.

\begin{theorem}[The support of every pair pole]\label{gf:thm:gf-nested-poles}
The recursive substitution defines an injective algebra map
\(\Phi_T:L_n\to\mathscr L_T\).
For the block \(S_e\) below an internal edge, define
\[
 N_0=\sum_{a<b}N_{ab},\qquad
 N_e=\sum_{\substack{a<b\\a,b\in S_e}}N_{ab}.
\]
Then
\begin{equation}\label{gf:eq:gf-pole-bound}
 \Phi_T(F_NL_n)\subseteq
 \rho^{-N_0}\prod_e t_e^{-N_e}\,
 Q_T[[t_e:e\in E(T)]].
\end{equation}
For \(1/D_N\), every displayed normal exponent is exact.
After removing these monomials, restriction to all \(t_e=0\)
is independent of the order of restriction.
Successive collision expansions along a fixed tree agree in
\(\mathscr L_T\) when they use these same relative scales.
\end{theorem}

\begin{proof}
Let \(v\) be the lowest common ancestor of \(a<b\), and
let \(c<d\) be their distinct children at \(v\).
Subtract the two recursive coordinate expressions.
Their common ancestral positions cancel.  Their common scales factor,
giving the exact identity
\begin{equation}\label{gf:eq:gf-pair-factor}
 \Phi_T(\delta_{ab})=
 \rho\prod_{e:\ a,b\in S_e}t_e\ U_{ab}(t),
 \qquad U_{ab}(0)=u_{v,d}-u_{v,c}.
\end{equation}
The expression \(U_{ab}\) is polynomial in the normal variables.
Its constant term is a unit of \(Q_T\).
It therefore has a unique inverse in the simultaneous formal
power-series ring.  Multiplying \eqref{gf:eq:gf-pair-factor} over
all pairs gives exactly the exponents in
\eqref{gf:eq:gf-pole-bound}.  Polynomial numerators again contribute
no negative powers.  For numerator \(1\), the remaining factor
is a unit even when all normal variables vanish.  The exponents
are consequently sharp.

The insertion coordinates are birational on their ordered open
domain.  To recover a vertex scale, take the separation between
the first leaves in its first and last children.  Divide it by
the analogous parent separation to recover \(t_e\).
The root separation recovers \(\rho\), and the first leaf
recovers \(b\).  Ratios of the other child anchors recover
all \(u_{v,j}\).  This gives the rational inverse and proves
injectivity, as in Theorem~\ref{gf:thm:gf-single-poles}.

After removing the specified negative normal monomials, every
expression is a simultaneous power series.  Setting any subset
of its variables to zero is defined and these evaluations commute.
Recursive insertion uses the same affine formulas whichever
compatible edges are processed first.  Each pair denominator has
the unique formal inverse of the unit \(U_{ab}\).
Uniqueness proves equality of the successive expansions in the
specified simultaneous ring.  It does not identify arbitrary
iterated Laurent fields with that ring.
\end{proof}

The outer scale retains the total collision as well.  In the
translation and scale compactification, it gives the face
\(\mathbb R\times\{\rho=0\}\times K_n\), locally up to a
positive change of scale.  The exponent there is \(N_0\).
The normalized associahedron omits this face because its diameter
is fixed.  Thus the pair \((\rho,K_n)\) records the total
collision and every nested relative collision separately.

\begin{example}[A five-point nested pole]\label{gf:ex:gf-nested}
Take the nested blocks \(234\supset23\).  The root has
children \((1,234,5)\), with coordinates \((0,a,1)\).
The other vertices are binary.  With scales \(t,s\), the positions are
\[
 (z_1,z_2,z_3,z_4,z_5)=
 \bigl(b,\ b+\rho a,\ b+\rho(a+ts),\ b+\rho(a+t),\ b+\rho\bigr).
\]
The open edge stratum has \(0<a<1\).  Its two boundary normals
are \(t,s\).  For
\[
 f=\frac1{\delta_{12}\delta_{23}^2\delta_{34}\delta_{45}},
\]
direct substitution gives
\begin{equation}\label{gf:eq:gf-example}
 \Phi_T(f)=\frac{\rho^{-5}t^{-3}s^{-2}}
                  {a(1-s)(1-a-t)}.
\end{equation}
The formula \(N_{234}=3\), \(N_{23}=2\), \(N_0=5\)
agrees with each exact exponent.
The coefficient after removing those monomials and setting
\(s=t=0\) is \(1/[a(1-a)]\), in either order.
The coefficient of \(t^ms^\ell\) in the remaining power series is
\[
 \frac1{a(1-a)^{m+1}},\qquad m,\ell\geq0.
\]
The growing pole at \(a=1\) is a tangential pole toward the
endpoint where the outer clusters \(234\) and \(5\) collide.
\end{example}

\subsubsection{The coefficient chain map at a collision}

The leading coefficient map extends to logarithmic differential forms.
Its target differential records the normal pole order.
Let \(\Omega^*_{L_n/\mathbb Q}\) have cohomological grading and
the algebraic de Rham differential \(d\).
Inside this complex, define
\[
 \Omega^*_{n,N}=
 F_NL_n\cdot
 \bigwedge\!\langle dz_1,d\log\delta_{ab}:a<b\rangle_{R_n}.
\]
This denotes the image in the actual de Rham algebra, so all
relations among these generators are imposed.
It is a subcomplex.  Indeed,
\(d(P/D_N)=dP/D_N-(P/D_N)\sum N_{ab}d\log\delta_{ab}\),
and
\(dz_j=dz_1+\delta_{1j}d\log\delta_{1j}\) for \(j>1\).
The union over \(N\) is the full rational de Rham complex.

In the target, introduce closed degree-one symbols
\(\lambda_0,\lambda_e\), representing \(d\log\rho,d\log t_e\).
Let \(\mathscr D_T^*\) be the exterior algebra over
\(\mathscr L_T\) on \(db\), all intermediate \(du_{v,j}\),
and these symbols.  Its continuous differential on coefficients is
\[
 d_Tf=(\partial_bf)db+
 \sum_{v,j}(\partial_{u_{v,j}}f)du_{v,j}
 +\rho(\partial_\rho f)\lambda_0
 +\sum_e t_e(\partial_{t_e}f)\lambda_e.
\]
It kills each displayed degree-one generator.
The commuting derivatives imply \(d_T^2=0\).
All formal derivatives are taken coefficientwise.
The simultaneous normal lower bounds make them well-defined.

Substitution and differentiation give the cochain algebra map
\begin{equation}\label{gf:eq:gf-dr-map}
 \Phi_T^{\mathrm{dR}}:
 (\Omega^*_{L_n/\mathbb Q},d)
 \longrightarrow(\mathscr D_T^*,d_T).
\end{equation}
Its values on collision forms are explicit:
\begin{equation}\label{gf:eq:gf-dlog}
 \Phi_T^{\mathrm{dR}}(d\log\delta_{ab})=
 \lambda_0+\sum_{e:\ a,b\in S_e}\lambda_e+d_T\log U_{ab}.
\end{equation}
Here \(d_T\log U\) means \(U^{-1}d_TU\), so no logarithm
function or branch is chosen.
Also \(\Phi_T^{\mathrm{dR}}(dz_1)=db\).

Write \(A_T\) for \(Q_T\) with the generator \(\rho\) omitted,
so \(Q_T=A_T[\rho]\).  The leading target is the graded module
\[
 \mathscr R_T^*=\Omega^*_{A_T[\rho,\rho^{-1}]/\mathbb Q}
                    \otimes\bigwedge\langle\lambda_e:e\in E(T)\rangle.
\]
Within the de Rham factor, \(\lambda_0=d\rho/\rho\).
Let \(d_B\) be its de Rham differential, extended by
\(d_B\lambda_e=0\), and put
\[
 \omega_N=N_0\lambda_0+\sum_eN_e\lambda_e,
 \qquad d_N\alpha=d_B\alpha-\omega_N\wedge\alpha.
\]
The closed one-form \(\omega_N\) has square zero, so
\(d_N^2=0\).  For nonzero \(N\), this is a complex with a
twisted differential, rather than a de Rham algebra differential.

\begin{theorem}[Leading logarithmic fusion]\label{gf:thm:gf-leading-chain}
The formula
\begin{equation}\label{gf:eq:gf-leading-chain}
 \mathcal R_{T,N}(\alpha)=
 \left.
 \rho^{N_0}\prod_et_e^{N_e}\,
 \Phi_T^{\mathrm{dR}}(\alpha)\right|_{t_e=0\text{ for all }e}
\end{equation}
defines a cochain map
\[
 \mathcal R_{T,N}:(\Omega^*_{n,N},d)
       \longrightarrow(\mathscr R_T^*,d_N).
\]
Evaluation in \eqref{gf:eq:gf-leading-chain} sets the coefficient
variables \(t_e\) to zero and retains the formal symbols \(\lambda_e\).
Restriction along several relative normal variables commutes.
For \(\alpha\in\Omega^*_{n,N}\) and
\(\beta\in\Omega^*_{n,M}\),
\[
 \mathcal R_{T,N+M}(\alpha\wedge\beta)
     =\mathcal R_{T,N}(\alpha)\wedge\mathcal R_{T,M}(\beta).
\]
\end{theorem}

\begin{proof}
The universal rule \(df\mapsto d_T\Phi_T(f)\) defines
\eqref{gf:eq:gf-dr-map} and makes it commute with differentials.
Differentiating \eqref{gf:eq:gf-pair-factor} proves
\eqref{gf:eq:gf-dlog}.  The coefficients of \(d_T\log U_{ab}\)
are power series with no negative normal powers.
Thus every form in \(\Omega^*_{n,N}\) obeys the same normal
bound as its coefficient in Theorem~\ref{gf:thm:gf-nested-poles}.
The right side of \eqref{gf:eq:gf-leading-chain} is therefore defined.

Put \(w_N=\rho^{N_0}\prod_et_e^{N_e}\).
The product rule gives, with no degree sign on \(w_N\),
\[
 d_T(w_N\Phi_T^{\mathrm{dR}}\alpha)
  =\omega_N\wedge(w_N\Phi_T^{\mathrm{dR}}\alpha)
       +w_N\Phi_T^{\mathrm{dR}}(d\alpha).
\]
The coefficient of normal degree zero in \(t_e\partial_{t_e}f\)
vanishes for a power series \(f\).
Hence coefficient evaluation intertwines \(d_T\) with \(d_B\)
on these weighted forms.  The displayed identity becomes
\(d_N\mathcal R_{T,N}(\alpha)=\mathcal R_{T,N}(d\alpha)\).
Commutativity of the coefficient evaluations proves the restriction
assertion.  Finally, \(w_{N+M}=w_Nw_M\) and evaluation is
multiplicative, which proves the last formula.
\end{proof}

For the function in Example~\ref{gf:ex:gf-nested}, write
\(h=1/[a(1-a)]\).  Denote the two relative normal symbols by
\(\lambda_t,\lambda_s\).  The leading maps give
\begin{gather*}
 \mathcal R_{T,N}(df)=d_Bh-(5\lambda_0+3\lambda_t+2\lambda_s)h,\\
 \mathcal R_{T,0}(d\log\delta_{23})=\lambda_0+\lambda_t+\lambda_s.
\end{gather*}
Using the ordinary differential \(d_B\) alone would omit all
three normal contributions.
The leading target keeps these normal directions, even when their
coefficient variables have been set to zero.

\begin{proposition}[Nonzero normal modes]\label{gf:prop:gf-normal-mode}
If \(N_e>0\) for some relative collision edge, then
\((\mathscr R_T^*,d_N)\) is contractible over \(\mathbb Q\).
In particular, leading logarithmic fusion does not identify this
complex with a coefficient interval that retains generic pole classes.
\end{proposition}

\begin{proof}
Let \(\iota_e\) be exterior contraction with \(\lambda_e\).
It kills base forms and the other normal symbols, and sends
\(\lambda_e\) to \(1\).  The tensor signs give
\(d_B\iota_e+\iota_ed_B=0\), while
\(\omega_N\wedge\iota_e+\iota_e(\omega_N\wedge-) = N_e\).
Thus \(h=-N_e^{-1}\iota_e\) satisfies
\(d_Nh+hd_N=1\).

For comparison, over \(R=\mathbb Q[\delta]\), the coefficient
interval with degree-zero module \(Re_0\oplus Re_1\), degree-one
module \(R[\delta^{-1}]\theta\), and differential
\(\partial e_0=-\theta\), \(\partial e_1=\theta\), has
\[
 H^1=R[\delta^{-1}]/R,
 \qquad [\theta/\delta]\ne0.
\]
Its degree-zero module does not contain \(e_1/\delta\).
The logarithmic normal complex instead admits
\(d(t^{-1})=-t^{-1}d\log t\).
These are different complexes with different boundary behavior.
\end{proof}

The relative-edge hypothesis is essential.  For the corolla there
are no relative normal symbols, and
\(d_N=d_B-N_0\,d\log\rho\wedge\).  Multiplication by the
invertible function \(\rho^{N_0}\) gives the chain isomorphism
\[
 (\mathscr R_T^*,d_B)\longrightarrow(\mathscr R_T^*,d_N),
 \qquad \alpha\longmapsto\rho^{N_0}\alpha,
\]
since \(d_N(\rho^{N_0}\alpha)=\rho^{N_0}d_B\alpha\).
Thus a positive dilation exponent alone does not imply contraction.

\begin{proposition}[A necessary limit on the coefficient target]
\label{gf:prop:gf-no-uniform-tangent}
Even a single rational pair pole can have no uniform tangential
pole bound across its formal normal coefficients.
The map \(\Phi_T\) cannot in general take values in the union
of coefficientwise products with one fixed tangential pole tuple.
\end{proposition}

\begin{proof}
The expansion
\[
 (A+t)^{-1}=\sum_{m\geq0}(-1)^mA^{-m-1}t^m
\]
has exact pole order \(m+1\) at the irreducible divisor \(A=0\).
No finite bound contains all its coefficients.
The factor \((1-a-t)^{-1}\) in
Example~\ref{gf:ex:gf-nested} realizes this phenomenon in an actual
five-point collision chart.  The normal lower bounds remain finite,
as \eqref{gf:eq:gf-example} shows.
\end{proof}

The restriction to the open stratum in \(Q_T\) is therefore
substantive.  To approach its boundary, one introduces the further
relative scales of the refined tree and applies
\eqref{gf:eq:gf-pair-factor} there.  Bare evaluation of a pole on
a closed face is not a defined fusion map.

\subsection{An operation on chains and the representation problem}
\label{gf:sec:gf-action}

There is a concrete chain algebra in which the operation
\(m_5\) is nonzero.  Put \(K_1=\{\mathrm{pt}\}\), with its
point as the insertion identity, and define
\[
 \mathcal A=\bigoplus_{r\geq1}C_*(K_r;\mathbb Z).
\]
Repeated collision insertion gives maps
\[
 C_*(K_n;\mathbb Z)\otimes
 \bigotimes_{j=1}^nC_*(K_{r_j};\mathbb Z)
 \longrightarrow C_*(K_{r_1+\cdots+r_n};\mathbb Z)
 \subseteq\mathcal A.
\]
Use the tensor differential and the orientation signs for every
permutation of factors.  Lemma~\ref{gf:lem:gf-grafting} proves their
composition identities.  These maps define an action of the
cellular collision operad on \(\mathcal A\).
If \(x\) is the generator of \(C_0(K_1;\mathbb Z)\), then
\[
 m_5(x,x,x,x,x)=m_5\in C_3(K_5;\mathbb Z)\subseteq\mathcal A.
\]
This is a nonzero chain.  Its boundary is the nonzero signed sum
of the nine distinct two-cells in \eqref{gf:eq:gf-nine}.
The cancellation occurs only after taking the next boundary.

For a prescribed cochain complex \((V,d_V)\), the required action
has a different target:
\begin{equation}\label{gf:eq:gf-end-action}
 C_d(K_n;\mathbb Q)\longrightarrow
 \operatorname{Hom}^{-d}(V^{\otimes n},V).
\end{equation}
Here
\(d_{\mathrm{Hom}}F=d_VF-(-1)^{|F|}F d_{V^{\otimes n}}\),
and tensor differentials and compositions use the Koszul convention.
An action means maps \eqref{gf:eq:gf-end-action} that respect this
differential and every collision insertion.
Such an action sends \(m_n\) to an operation \(M_n\) of
degree \(2-n\).  Theorem~\ref{gf:thm:gf-boundary} then gives
\eqref{gf:eq:gf-nine} with \(d_{\mathrm{Hom}}M_5\) on the left
and the corresponding nine composites on the right.
This consequence applies only when that action is specified.

\begin{problem}[Collision action on polynomial states]
\label{gf:prob:gf-native-action}
Let \(V\) be the one-species free state algebra
\eqref{fs:eq:state-algebra} over \(\mathbb Q\).
Its pair contraction coefficients \(K_{rs}(z_a-z_b)\)
have normal exponents prescribed by
Theorem~\ref{gf:thm:gf-nested-poles}.  Construct chain operations
over the collision strata that realize these contractions, their
translation maps, and their interval coefficients.
The construction must identify the state carrier over each stratum
and give comparison maps on every nested and disjoint collision.
Its five-input boundary must equal the nine composites in
\eqref{gf:eq:gf-nine}, after transport into one common target complex.
For mode-module bars, it must also preserve the two module actions
and every balanced cut.
\end{problem}

The chain maps \(\Gamma\) and the rational maps \(\Phi\)
provide different parts of this problem.  The first use integral
cellular chains on closed geometric strata.  The second use rational
germs in formal normal neighborhoods of open strata.
Neither construction identifies interval overlap classes with chains
on those strata.  A representation must supply that comparison.
The nonuniform tangential poles in
Proposition~\ref{gf:prop:gf-no-uniform-tangent} also require a stated
topology before completed state or bar inputs are admitted.
No interacting factorization algebra or completed bar equivalence
is asserted by these finite geometric constructions.
